\documentclass[11pt]{article}
\usepackage[margin=1in]{geometry}
\usepackage{amsmath,amssymb,amsthm,hyperref}
\newtheorem{theorem}{Theorem}
\newtheorem{lemma}{Lemma}
\newcommand{\F}{\mathbb F_2}
\title{Inverting self-complementary symmetric graphs}
\author{Principia Math In Progress}
\date{29 September 2026}
\begin{document}
\maketitle

\paragraph{Task 1 (original problem).}
MathDB \href{https://mathdb.com/p/380331}{380331} records
Genest's Conjecture 1.7.13 \cite[p.~39]{Genest}.
A finite graph is \emph{symmetric} here if its automorphism group is
transitive on vertices and on edges. For a self-complementary symmetric
graph \(G\) of positive order \(n\), the conjecture asserts
\[
 G\text{ is invertible}\quad\Longleftrightarrow\quad
 n=1\text{ or }n\equiv5\pmod8,
 \qquad G^{-1}=\overline G\text{ whenever it exists}.
\]
The inverse in this statement is Genest's graph operation, not an
assumption about an adjacency matrix. We first give its definitions and
prove the required connection.

\paragraph{The source operation.}
Give every vertex a white or black color. Local complementation at \(u\)
toggles the edges between distinct neighbors of \(u\). If \(u\) is white,
it also reverses all its neighbors' colors; if \(u\) is black, it changes
no color. The vertex \(u\) keeps its own color.
A \emph{parity word} is built from allowed steps of two kinds:
a local complementation at a white vertex \(u\), or the three successive
local complementations \(uvu\), abbreviated \([uv]\), at adjacent
black vertices \(u,v\).
A reduced parity word is composed of singleton and pair blocks with
disjoint vertex supports. The graph is invertible when such a word
uses every vertex, and its inverse is the resulting graph.
These are Definitions 1.3.9, 1.6.4, and 1.7.1 in \cite{Genest};
the source proves independence of the choice of reduced word.
For an uncolored graph, its \emph{natural coloring} makes even-degree
vertices white and odd-degree vertices black.

\begin{lemma}\label{matrix}
For a bicolored graph \(H\), form the symmetric matrix \(M_H\) over
\(\F\) whose off-diagonal entries are its adjacencies, and whose
diagonal is one at white vertices and zero at black vertices.
Then \(H\) is invertible in the source sense if and only if \(M_H\)
is nonsingular. In that case
\[
 M_{H^{-1}}=M_H^{-1}.
\]
\end{lemma}
\begin{proof}
Partition a symmetric matrix as
\[
 M=\begin{pmatrix}P&Q\\Q^t&R\end{pmatrix},
\]
where \(P\) is nonsingular. Solving the first block of the equation
\(y=Mx\) for \(x_S\), where \(S\) indexes \(P\), gives
\begin{equation}\label{pivot}
 \begin{pmatrix}x_S\\y_{S^c}\end{pmatrix}
 =
 \begin{pmatrix}
 P^{-1}&P^{-1}Q\\
 Q^tP^{-1}&R+Q^tP^{-1}Q
 \end{pmatrix}
 \begin{pmatrix}y_S\\x_{S^c}\end{pmatrix}.
\end{equation}
All signs are plus in characteristic two.
This operation is a principal pivot; its displayed formula is all
we need.

For a white vertex, \(P=(1)\). Formula \eqref{pivot} preserves the
incident edges and pivot diagonal, toggles all edges between its
neighbors, and toggles exactly their diagonal entries. It is exactly
the source's white-vertex move.

For adjacent black vertices \(u,v\),
\[
 P=P^{-1}=\begin{pmatrix}0&1\\1&0\end{pmatrix},\qquad
 Q=\begin{pmatrix}c^t\\d^t\end{pmatrix}.
\]
The transformed matrix is
\begin{equation}\label{pair}
 \begin{pmatrix}
 0&1&d^t\\
 1&0&c^t\\
 d&c&R+cd^t+dc^t
 \end{pmatrix}.
\end{equation}
The source move \(uvu\) has this same effect.
For an outside vertex whose original adjacencies to \(u,v\) are
\((c_i,d_i)\), the three moves change them to \((d_i,c_i)\).
For outside vertices \(i,j\), the three edge toggles sum to
\[
 c_ic_j+(c_i+d_i)(c_j+d_j)+d_id_j
 =c_id_j+d_ic_j.
\]
All three pivots are at black vertices, so colors are unchanged,
also agreeing with the zero diagonal of \(cd^t+dc^t\).

It remains to check existence of a word covering all vertices.
Block elimination gives
\[
 \det M=\det P\,\det(R+Q^tP^{-1}Q).
\]
If \(M\) is nonsingular, either some diagonal entry is one, allowing
a white singleton, or all diagonal entries are zero and some
off-diagonal entry is one, allowing an adjacent black pair.
The remaining Schur complement is symmetric and nonsingular.
Repeat on the as yet unused vertices. The moves are allowed in
the current graph and have disjoint supports, so they give a reduced
parity word covering every vertex.

Conversely, each block of any such full reduced word is a nonsingular
singleton or pair pivot. On the unused vertices, the successive
matrices are exactly the successive Schur complements. The determinant
factorization, ending in the empty determinant one, shows that the
original matrix is nonsingular.

Finally \eqref{pivot} interchanges the roles of \(x_i,y_i\) precisely
on the pivot set. Successive disjoint pivots covering all indices
therefore transform \(y=Mx\) into \(x=M^{-1}y\). Their final matrix is
\(M^{-1}\), proving both the formula and its independence of the word.
\end{proof}

\begin{lemma}\label{parameters}
Let \(G\) be a self-complementary symmetric graph of order \(n>1\).
Then \(n\equiv1\pmod4\), and \(G\) is strongly regular with parameters
\[
 k=\frac{n-1}{2},\qquad \lambda=\frac{n-5}{4},\qquad
 \mu=\frac{n-1}{4}.
\]
\end{lemma}
\begin{proof}
Vertex transitivity makes the degree constant, and self-complementarity
gives \(k=n-1-k\). Thus \(n\) is odd; since the number of edges is
\(n(n-1)/4\), we get \(n\equiv1\pmod4\).
For \(n>1\) this implies \(n\ge5\), so edges and nonedges both exist.

Edge transitivity makes the number \(\lambda\) of common neighbors
constant on adjacent pairs. An isomorphism to the complement
conjugates edge transitivity to transitivity on nonadjacent pairs,
so their common-neighbor count \(\mu\) is constant too.
In the complement, an adjacent pair has \(n-2-2k+\mu=\mu-1\)
common neighbors. Hence \(\lambda=\mu-1\).
Counting length-two walks from a fixed vertex to its nonneighbors
gives
\[
 k(k-\lambda-1)=(n-k-1)\mu=k\mu.
\]
Since \(k>0\), we obtain \(k-\lambda-1=\mu\). Combining these equations
gives \(k=2\mu\), and the displayed parameters follow.
\end{proof}

\begin{theorem}
Genest's invertibility conjecture holds in full.
\end{theorem}
\begin{proof}
For \(n=1\), the sole vertex has degree zero, is white, and its local
complementation leaves the graph unchanged. It is its own inverse and
its own complement.

Suppose \(n>1\). Let \(A\) be the adjacency matrix of \(G\), let \(I\)
be the identity matrix, and let \(J\) be the all-ones matrix.
Lemma~\ref{parameters} gives the integer identity
\begin{equation}\label{srg}
 A^2=(k-\mu)I+(\lambda-\mu)A+\mu J.
\end{equation}
Because \(k=(n-1)/2\) is even, the natural coloring is all white;
therefore the matrix in Lemma~\ref{matrix} is \(M=I+A\).
In what follows, all matrix operations are over \(\F\).
Also \(AJ=kJ=0\).

If \(n\equiv1\pmod8\), then \(\mu\) and \(k-\mu\) are even, while
\(\lambda-\mu=-1\) is odd. Equation \eqref{srg} becomes \(A^2=A\).
Thus
\[
 (I+A)A=0.
\]
There is an edge, so \(A\ne0\), and \(I+A\) has a nonzero kernel.
Lemma~\ref{matrix} shows that \(G\) is not invertible.

If \(n\equiv5\pmod8\), then \(\mu=k-\mu\) is odd and
\eqref{srg} becomes \(A^2=I+A+J\). Hence
\[
 (I+A)(A+J)=A+A^2+J+AJ=I,
 \qquad (I+A)^{-1}=A+J.
\]
The last matrix has diagonal one and off-diagonal entries
\(1+A_{ij}\). It therefore encodes \(\overline G\) with all vertices
white. The complement has the same even degree \(k\), so this is its
natural coloring. Lemma~\ref{matrix} proves
\(G^{-1}=\overline G\), with equality on the original labeled vertex set.
These two residue classes exhaust the possibilities.
\end{proof}

\paragraph{Scope and verification.}
The proof establishes both directions and identifies the full inverse,
including its coloring. It applies to every self-complementary symmetric
graph and does not require the classification into Paley and Peisert
graphs and the exceptional graph. The pivot correspondence is standard
linear algebra, proved here to make the source's convention explicit.
The accompanying exact verifier independently simulates the source
graph moves for every bicolored graph on at most four vertices and
compares feasible supports with principal minors. It also checks the
theorem on larger Paley graphs over prime and extension fields.
Finite checks supplement, rather than replace, the general proof.

\paragraph{Task 2 (proposed follow-up).}
Determine the minimum number of black-edge pair moves among reduced
words realizing the inverse of each self-complementary symmetric graph.

\paragraph{Task 3 (proposed follow-up).}
Determine the number and size distribution of complementation sets
in Paley and Peisert graphs, equivalently the nonsingular principal
submatrices of their color-augmented adjacency matrices.

\paragraph{Task 4 (proposed follow-up).}
Classify invertibility and inverse graphs for self-complementary
vertex-transitive graphs without the edge-transitivity assumption.

\paragraph{Task 5 (proposed follow-up).}
Construct full reduced inversion words directly from finite-field
coordinates of Paley graphs, with a formula that avoids matrix elimination.

\begin{thebibliography}{1}
\bibitem{Genest}
F. Genest, \emph{Graphes eul\'eriens et compl\'ementarit\'e locale}
(English arXiv title: \emph{Eulerian graphs and local complementation}),
Ph.D. thesis, Universit\'e de Montr\'eal; arXiv:math/0701421v1 (2007),
103 pages. Chapter 1, Definitions 1.3.9, 1.6.4, 1.7.1 and
Conjecture 1.7.13.
\url{https://arxiv.org/abs/math/0701421}.
\end{thebibliography}
\end{document}
